\begin{figure}[ht]
\centering
\begin{tikzpicture}[x=0.022cm, y=0.55cm, font=\small]

% Horizontal axis: calendar year from 1650 to 2030
\draw[->, thick] (1640, 0) -- (2050, 0) node[right] {year};
\foreach \x in {1700, 1750, 1800, 1850, 1900, 1950, 2000} {
  \draw[thick] (\x, -0.15) -- (\x, 0.15);
  \node[anchor=north] at (\x, -0.2) {\x};
}

% Vertical axis: log10 of fractional uncertainty (negative)
\draw[->, thick] (1640, 0) -- (1640, 11) node[above]
  {$-\log_{10}(\text{fractional uncertainty})$};
\foreach \y/\lbl in {2/2, 4/4, 6/6, 8/8, 10/10, 12/12, 14/14, 16/16, 18/18} {
  \draw[thick] (1635, \y) -- (1645, \y);
  \node[anchor=east] at (1632, \y) {\lbl};
}

% Eight epochal clocks: (year, -log10 fractional uncertainty), labelled
% Huygens 1656 pendulum: 10 s/day = 10/86400 = 1.16e-4, -log10 = 3.9 -> we
% plot 4 (approx)
\fill[engblack] (1656, 4) circle (3pt);
\node[anchor=west, font=\scriptsize] at (1660, 4)
  {Huygens pendulum (1656)};

% Harrison H4 1761 marine chronometer: 1 s/day = 1.16e-5 -> 4.9
\fill[engblack] (1761, 5) circle (3pt);
\node[anchor=west, font=\scriptsize] at (1765, 5)
  {Harrison H4 (1761)};

% Shortt-Synchronome 1921 free-pendulum: 0.01 s/day = 1.16e-7 -> 6.9
\fill[engblack] (1921, 7) circle (3pt);
\node[anchor=east, font=\scriptsize] at (1916, 7)
  {Shortt-Synchronome (1921)};

% Marrison quartz 1927: 1 ms/day = 1.16e-8 -> 7.9
\fill[engblue] (1927, 8) circle (3pt);
\node[anchor=west, font=\scriptsize] at (1931, 8)
  {Marrison quartz (1927)};

% Essen-Parry caesium beam 1955: 1e-10 (early); -log10 = 10
\fill[engred] (1955, 10) circle (3pt);
\node[anchor=east, font=\scriptsize] at (1950, 10)
  {Essen-Parry Cs (1955)};

% NIST-F1 caesium fountain 1999: 1e-15 -> 15
\fill[engred] (1999, 15) circle (3pt);
\node[anchor=west, font=\scriptsize] at (2003, 15)
  {NIST-F1 Cs fountain (1999)};

% Strontium optical lattice 2010+ (use 2014): 1e-18 -> 18
\fill[engblue] (2014, 18) circle (3pt);
\node[anchor=east, font=\scriptsize] at (2009, 18)
  {Sr optical lattice (2014)};

% CSAC 2011 (chip-scale): 1e-11 -> 11
\fill[engblue] (2011, 11) circle (3pt);
\node[anchor=west, font=\scriptsize] at (2015, 11)
  {CSAC chip-scale (2011)};

% Connecting trend line (eyeball, not regression)
\draw[dashed, gray]
  (1656, 4) -- (1761, 5) -- (1921, 7) -- (1927, 8)
  -- (1955, 10) -- (1999, 15) -- (2014, 18);

% Annotate eras
\node[engblack, font=\footnotesize, anchor=south]
  at (1750, 0.4) {mechanical};
\node[engblue, font=\footnotesize, anchor=south]
  at (1930, 0.4) {quartz};
\node[engred, font=\footnotesize, anchor=south]
  at (1975, 0.4) {microwave atomic};
\node[engblue, font=\footnotesize, anchor=south]
  at (2015, 0.4) {optical};

\end{tikzpicture}
\caption[Time-precision arc 1656-2020]{Approximate fractional
frequency uncertainty $\sigma_{y}$ of the working time standard at
each epoch, plotted on a $-\log_{10}$ axis against calendar year.
Each marker represents a step change in achievable precision: the
Huygens pendulum (1656) at $\sim 10^{-4}$, Harrison H4 (1761) at
$\sim 10^{-5}$, the Shortt-Synchronome free-pendulum (1921) at
$\sim 10^{-7}$, Marrison's quartz crystal clock (1927) at
$\sim 10^{-8}$, Essen and Parry's caesium beam (1955) at
$\sim 10^{-10}$, the NIST-F1 caesium fountain (1999) at
$\sim 10^{-15}$, and the JILA/NIST strontium optical lattice
clock (2014) at $\sim 10^{-18}$. The chip-scale atomic clock
(CSAC, $\sim 2011$) sits at $\sim 10^{-11}$, trading ultimate
accuracy for portability and power. Approximate trend: roughly
one decade of fractional precision per generation since 1650.
Sources: \cite{paper:ch06t3-marrison-1948,
paper:ch06t3-essen-parry-1955, web:nist-cesium-fountain,
web:ch06t3-nist-strontium-2024, web:ch06t3-microsemi-csac-2024,
hist:ch06t3-sobel-longitude-1995}.}
\label{fig:vol01:ch06:time-precision-arc}
\end{figure}
